\subsection{Stärken}
\label{subsec:staerken}

Die gemeinsame Baseline verbindet zentrale Pflege mit klaren Grenzen zwischen
Frontend, Backend, Desktop, Persistenz und Tooling. Das geteilte Frontend
adressiert A-02 und A-05, ohne jedes Profil mit allen Komponenten zu belasten.
Diese Trennung unterstützt eigene Tests und produktspezifische Erweiterungen,
ohne die gemeinsame Präsentationsschicht zu duplizieren.

Fünf Profile bilden wiederkehrende Plattformzuschnitte ab. Abhängigkeiten
verhindern PostgreSQL ohne Backend. Die Abnahme aller Basisprofile sowie der
zulässigen und unzulässigen PostgreSQL-Kombinationen stützt dieses Modell für
den begrenzten aktuellen Katalog, nicht für beliebig viele künftige
Capabilities \parencite{kleiveist2026acceptance}.

\texttt{control.py} vereinheitlicht den Einstieg, während Fachwerkzeuge
sichtbar bleiben. Manifest, Konfigurationsvertrag, Tests und Buildpfade erhöhen
die Nachvollziehbarkeit. Die byte-identische Generierung zweier
\texttt{desktop-cloud + postgres}-Projekte liefert dafür einen konkreten, aber
auf eine Kombination und Umgebung begrenzten Nachweis
\parencite{kleiveist2026acceptance}.
Profilabhängiges Überspringen nicht vorhandener Schichten hält die gemeinsame
Befehlsoberfläche dabei an den tatsächlichen Projektumfang gebunden.

Tabelle~\ref{tab:strengths-limitations} stellt Stärken und Geltungsgrenzen
gegenüber.

\begin{table}[htbp]
  \centering
  \scriptsize
  \renewcommand{\arraystretch}{1.08}
  \caption{Evidenzbasierte Gegenüberstellung von Stärken und Grenzen}
  \label{tab:strengths-limitations}
  \begin{tabularx}{\textwidth}{@{}>{\raggedright\arraybackslash}p{0.14\textwidth}>{\raggedright\arraybackslash}p{0.27\textwidth}>{\raggedright\arraybackslash}p{0.27\textwidth}>{\raggedright\arraybackslash}X@{}}
    \toprule
    \textbf{Bereich} & \textbf{Stärke und Evidenz} & \textbf{Grenze und Evidenz} & \textbf{Bedeutung} \\
    \midrule
    Architektur & Getrennte Frontend-, Backend-, Tauri- und Persistenzgrenzen;
      gemeinsames Frontend. & Shared Schemas werden nicht automatisch von
      Frontend und Backend konsumiert. & Wiederverwendung bei klarer
      Zuständigkeit; Vertragsänderungen bleiben manuell. \\
    Variabilität & Fünf deklarative Profile; Abhängigkeiten und drei zulässige
      PostgreSQL-Kombinationen sind abgenommen. & Nur eine auswählbare
      Capability; direkte CLI trennt \texttt{selectable} nicht strikt. &
      Kontrollierte Auswahl im aktuellen Katalog, begrenzte Aussage für
      Erweiterungen. \\
    Tooling & Gemeinsamer profilabhängiger Einstieg für Diagnose,
      Installation, Start, Test und Build. & Mehrere Fachwerkzeuge und native
      Voraussetzungen bleiben erforderlich. & Weniger unterschiedliche
      Bedienpfade, aber kein technologieunabhängiger Workflow. \\
    Qualität & Zentrale Policy, 13 Rule IDs, Exception- und Architekturtests;
      135 fokussierte Tests und grüner Master-Quality-Job. & Backendhaltige
      generierte Profile scheitern im Quality-Schritt; lokaler Gesamtlauf durch
      Rust-Systembibliotheken begrenzt; \texttt{CQ001} unterdrückbar. &
      Definierte Verstöße werden reproduzierbarer erkannt; allgemeine
      Codequalität oder Wartbarkeit ist nicht garantiert. \\
    Reproduzierbarkeit & Zwei gleich konfigurierte Projekte waren im
      dokumentierten Versuch byte-identisch. & Nachweis nur für eine
      Kombination und eine Untersuchungsumgebung. & Konkrete Evidenz ohne
      Anspruch auf vollständige plattformübergreifende Reproduzierbarkeit. \\
    Betrieb und Plattformen & Web- und Image-Builds, PostgreSQL-Integration
      sowie unsignierte Linux-/macOS-Pakete sind dokumentiert. & Compose-Start,
      produktiver Cloudbetrieb, Windows-Paket, Signing und Notarisierung
      bleiben offen. & Template-Reife ist nicht mit Produktionsreife
      gleichzusetzen. \\
    Wartung & Core, Profile, Tooling, Tests und Dokumentation liegen in einer
      Baseline. & Änderungen können mehrere Profile betreffen; Produktprojekte
      erhalten keine automatischen Updates. & Zentrale Pflege wird gebündelt,
      erfordert aber eine breite Regression und Transferprozesse. \\
    \bottomrule
  \end{tabularx}
  \smallskip
  {\footnotesize Quelle: Eigene Bewertung auf Grundlage der
  Kapitel~\ref{sec:architektur}--\ref{sec:qualitaetssicherung-verifikation}
  sowie des Abnahmeprotokolls und der aktuellen CI-Läufe
  \parencite{kleiveist2026acceptance,kleiveist2026ci-core-run,kleiveist2026ci-profiles-run,kleiveist2026ci-postgres-run,kleiveist2026ci-desktop-run,kleiveist2026quality-implementation,kleiveist2026governance-core-run,kleiveist2026governance-profiles-run}.
  \par}
\end{table}

